\section{Limitations and Future Work}\label{sec:limits}
\textbf{Synthetic data.} Every quantitative result describes behaviour on data generated by our own process, the situation against which the ``parable of Google Flu''~\cite{lazer} warns. Detection sensitivity depends on the size and shape of the injected outbreaks; forecast skill depends on how strongly the generator ties dengue to lagged weather, and the weather-model advantage in particular is built in by construction. Results on real DHIS2 data would be less favourable. The weather used for the reported results is synthetic (Open-Meteo could not be reached from the preparation environment); the fetch code is untested against the live API. Populations are approximate.

\textbf{Scale.} The reported runs use SF=0.05 on one virtual CPU; scaling beyond is extrapolated or must be measured on the team's laptop (SF=1, 25 million records). Spark was used in local mode; behaviour on a real cluster was not tested.

\textbf{Not implemented.} HDFS or an object store (ING-5), Delta Lake or Iceberg with time travel (ING-7), a k-anonymity check and differential privacy prototype (PRV-4, discussed only), the SEIR scenario module (AN-10), and the mobile or playback dashboard view (VIS-7). Weather in the forecast is not itself forecast: horizons up to four weeks use only observed weather four or more weeks old, which avoids the problem but limits the benefit.

\textbf{Methodological.} Forecast intervals ignore parameter uncertainty and so tend to be too narrow for series with little data; demand upper bounds sum weekly quantiles and are conservative; alert episodes and false-alert counts depend on the definition of an episode and a tolerance window around each true outbreak. The artefact guard handles one kind of artefact (back-fill after silence); a change of case definition would need a different guard. The usability study with non-specialist users is pending.

\textbf{Future work.} Nowcasting of the incomplete recent weeks from the reporting triangle (a moving-average completion factor would let horizon one use the last four weeks); an ensemble that adds a gradient-boosted model for accuracy while keeping the GLM as the audited reference; integer programmes with lead times and multiple periods; streaming alerts at daily resolution; and real data through a DHIS2 connector with the same bronze contract.
